\section*{Jak s materiálem pracovat}
\addcontentsline{toc}{section}{Jak s materiálem pracovat}
Tyto poznámky vycházejí z přiloženého PDF \emph{Statnice B-CS blok 3, otázky 15-23, v4.2}, zachovávají jeho právní korekce a doplňují je do podoby širších studijních zápisků. U prvních tří BSS otázek jsou navíc zapracované dodané starší poznámky: širší vymezení bezpečnostních studií, vztah bezpečnostních a strategických studií, sektory, pojmy hrozba/riziko a základy kyberválky. U právních otázek je hlavním východiskem ověřený aktuální stav, ne starší právo.

U každé otázky je nejdříve kostra odpovědi, potom rozbor přesně podle podotázek ze zadání. U zkoušky je dobré nemluvit jen seznamem pojmů, ale vždy ukázat vztah: \emph{co chráním, před kým nebo čím, jaký je právní/institucionální režim, jaký je praktický dopad a kde je typická zkoušková chyba}.

